\section{Contoh Terhitung}

\begin{example}
\textbf{Satu putaran jaringan Feistel.} Blok mainan 8 bit dibagi dua bagian 4 bit; fungsi putaran $f(R,K)=(R\oplus K)$ digeser siklik ke kiri 1 bit. Mula-mula $L_0=1010$, $R_0=0011$, dengan aturan
\[
L_i=R_{i-1},\qquad R_i=L_{i-1}\oplus f(R_{i-1},K_i).
\]
Putaran 1, $K_1=1100$: $f(0011,1100)=\text{geser-kiri}(0011\oplus1100)=\text{geser-kiri}(1111)=1111$, sehingga $L_1=0011$ dan $R_1=1010\oplus1111=0101$.\\
Putaran 2, $K_2=1010$: $f(0101,1010)=\text{geser-kiri}(1111)=1111$, sehingga $L_2=0101$ dan $R_2=0011\oplus1111=1100$.\\
Melanjutkan dengan $K_3=0110$ dan $K_4=1001$:
\[
\begin{array}{c|cc}
i & L_i & R_i\\
\hline
0 & 1010 & 0011\\
1 & 0011 & 0101\\
2 & 0101 & 1100\\
3 & 1100 & 0000\\
4 & 0000 & 1111
\end{array}
\]
Dekripsi memakai aturan yang sama dengan urutan kunci terbalik $(K_4,K_3,K_2,K_1)$ dan menukar peran $L,R$, tanpa perlu fungsi $f^{-1}$ yang terpisah.
\end{example}

\begin{example}
\textbf{Longsoran dan pembangkitan subkunci DES.}
(i) \textit{Efek longsoran.} Ubah satu bit plainteks di atas, $L_0:1010\to0010$. Setelah putaran 1, 2, 3, dan 4, jumlah bit keluaran $(L_i,R_i)$ yang berbeda dari semula berturut-turut adalah $1,2,3,3$ bit (dari 8 bit). Satu bit masukan makin menyebar seiring bertambahnya putaran; pada DES penuh (16 putaran, blok 64 bit) perubahan 1 bit menyebar ke sekitar setengah bit cipherteks.
(ii) \textit{Pembangkitan subkunci.} Kunci efektif 56 bit dipecah menjadi $C_0,D_0$ (masing-masing 28 bit). Untuk $C_0=1111000011110000111100001111$:
\[
\text{putaran 1, geser-kiri 1}\to C_1=1110000111100001111000011111,
\]
\[
\text{putaran 2, geser-kiri 1 lagi}\to C_2=1100001111000011110000111111,
\]
sesuai jadwal pergeseran putaran 1 dan 2 (masing-masing 1 bit). Hasil $C_i,D_i$ lalu dipilih 48 bit oleh PC-2 menjadi $K_i$.
(iii) \textit{Substitusi S-box.} S-box $S_1$ memetakan 6 bit ke 4 bit: bit terluar menentukan baris, empat bit tengah menentukan kolom. Untuk masukan $011011$: baris $01_2=1$, kolom $1101_2=13$, dan $S_1(1,13)=5$, jadi keluarannya $0101$.
\end{example}
